% GENERATED FILE. Edit canonical lessons, then run npm run generate:books.
% canonical-source-hash: fnv1a64:116ca691a16665aa
% canonical-lessons: UR-C129-pith, UR-C129-matha, UR-C129-gal, UR-C129-thori, UR-C129-kalai

\chapter{Back, Forehead, Cheek, Chin, Wrist}
\label{ch:ur-back-forehead-cheek-chin-wrist}
\hlchaptermodality{129}

\begin{chapteropening}
\textbf{By the end of this chapter:} \emph{I can say the back, the forehead, a cheek, the chin and a wrist.}

Complete the last lesson of chapter 129: I can say the back, the forehead, a cheek, the chin and a wrist.
\end{chapteropening}

\section[\texorpdfstring{\ur{پیٹھ} (pīṭh)}{پیٹھ (pīṭh)}]{\ur{پیٹھ} (pīṭh) — the back}

\label{lesson:UR-C129-pith}



\begin{quote}
Before the new one: say the Urdu for a bear, then the Urdu for the waist.
\end{quote}

\subsection*{You'll want to know: \ur{پیٹھ}}
\textbf{\ur{پیٹھ}} — \emph{pīṭh} — \textquotedblleft{}the back\textquotedblright{}.

\begin{cousinweb}
\textbf{\ur{پیٹھ}} (\emph{pīṭh}) is a tadbhava, a word inherited from Sanskrit and reshaped by sound change. It comes from \emph{pṛṣṭi}, \textquotedblleft{}vertebrae; back\textquotedblright{} (CDIAL 8370).
\end{cousinweb}

\begin{grammarlens}[title={What you've built}]
One more word for this chapter.
\end{grammarlens}

\subsection*{Your turn}
\begin{itemize}
\raggedright
  \item \emph{Say it:} \emph{pīṭh}
  \item \emph{Say it:} \emph{pīṭh}, once more
  \item \emph{Say it:} say \emph{pīṭh}
  \item \emph{Recall:} say the Urdu for a bear, then the Urdu for the waist, then say \emph{pīṭh} again
\end{itemize}

\begin{tcolorbox}[breakable,colback=teal!4,colframe=teal!35!black,title={Before you move on}]
What does \ur{پیٹھ} mean? (\textquotedblleft{}The back\textquotedblright{}.) Say it once more.
\end{tcolorbox}
\section[\texorpdfstring{\ur{ماتھا} (māthā)}{ماتھا (māthā)}]{\ur{ماتھا} (māthā) — the forehead}

\label{lesson:UR-C129-matha}



\begin{quote}
Before the new one: say the Urdu for the waist, then the Urdu for the back.
\end{quote}

\subsection*{You'll want to know: \ur{ماتھا}}
\textbf{\ur{ماتھا}} — \emph{māthā} — \textquotedblleft{}the forehead\textquotedblright{}.

\begin{cousinweb}
\textbf{\ur{ماتھا}} (\emph{māthā}) is a tadbhava, a word inherited from Sanskrit and reshaped by sound change. It comes from \emph{mastaka}, \textquotedblleft{}head, skull\textquotedblright{} (CDIAL 9926).
\end{cousinweb}

\begin{grammarlens}[title={What you've built}]
One more word for this chapter.
\end{grammarlens}

\subsection*{Your turn}
\begin{itemize}
\raggedright
  \item \emph{Say it:} \emph{māthā}
  \item \emph{Say it:} \emph{māthā}, once more
  \item \emph{Say it:} say \emph{māthā}
  \item \emph{Recall:} say the Urdu for the waist, then the Urdu for the back, then say \emph{māthā} again
\end{itemize}

\begin{tcolorbox}[breakable,colback=teal!4,colframe=teal!35!black,title={Before you move on}]
What does \ur{ماتھا} mean? (\textquotedblleft{}The forehead\textquotedblright{}.) Say it once more.
\end{tcolorbox}
\section[\texorpdfstring{\ur{گال} (gāl)}{گال (gāl)}]{\ur{گال} (gāl) — a cheek}

\label{lesson:UR-C129-gal}



\begin{quote}
Before the new one: say the Urdu for the back, then the Urdu for the forehead.
\end{quote}

\subsection*{You'll want to know: \ur{گال}}
\textbf{\ur{گال}} — \emph{gāl} — \textquotedblleft{}a cheek\textquotedblright{}.

\begin{cousinweb}
\textbf{\ur{گال}} (\emph{gāl}) is a tadbhava, a word inherited from Sanskrit and reshaped by sound change. It comes from \emph{galla}, \textquotedblleft{}cheek\textquotedblright{} (CDIAL 4089). CDIAL notes the view that it comes from Dravidian.
\end{cousinweb}

\begin{grammarlens}[title={What you've built}]
One more word for this chapter.
\end{grammarlens}

\subsection*{Your turn}
\begin{itemize}
\raggedright
  \item \emph{Say it:} \emph{gāl}
  \item \emph{Say it:} \emph{gāl}, once more
  \item \emph{Say it:} say \emph{gāl}
  \item \emph{Recall:} say the Urdu for the back, then the Urdu for the forehead, then say \emph{gāl} again
\end{itemize}

\begin{tcolorbox}[breakable,colback=teal!4,colframe=teal!35!black,title={Before you move on}]
What does \ur{گال} mean? (\textquotedblleft{}A cheek\textquotedblright{}.) Say it once more.
\end{tcolorbox}
\section[\texorpdfstring{\ur{ٹھوڑی} (ṭhoṛī)}{ٹھوڑی (ṭhoṛī)}]{\ur{ٹھوڑی} (ṭhoṛī) — the chin}

\label{lesson:UR-C129-thori}



\begin{quote}
Before the new one: say the Urdu for the forehead, then the Urdu for a cheek.
\end{quote}

\subsection*{You'll want to know: \ur{ٹھوڑی}}
\textbf{\ur{ٹھوڑی}} — \emph{ṭhoṛī} — \textquotedblleft{}the chin\textquotedblright{}.

\begin{cousinweb}
\textbf{\ur{ٹھوڑی}} (\emph{ṭhoṛī}) is a tadbhava, a word inherited from Sanskrit and reshaped by sound change. It comes from \emph{tuṇḍa}, \textquotedblleft{}beak, snout\textquotedblright{} (CDIAL 5853). CDIAL calls the word certainly non-Aryan and probably Munda.
\end{cousinweb}

\begin{grammarlens}[title={What you've built}]
One more word for this chapter.
\end{grammarlens}

\subsection*{Your turn}
\begin{itemize}
\raggedright
  \item \emph{Say it:} \emph{ṭhoṛī}
  \item \emph{Say it:} \emph{ṭhoṛī}, once more
  \item \emph{Say it:} say \emph{ṭhoṛī}
  \item \emph{Recall:} say the Urdu for the forehead, then the Urdu for a cheek, then say \emph{ṭhoṛī} again
\end{itemize}

\begin{tcolorbox}[breakable,colback=teal!4,colframe=teal!35!black,title={Before you move on}]
What does \ur{ٹھوڑی} mean? (\textquotedblleft{}The chin\textquotedblright{}.) Say it once more.
\end{tcolorbox}
\section[\texorpdfstring{\ur{کلائی} (kalāī)}{کلائی (kalāī)}]{\ur{کلائی} (kalāī) — a wrist}

\label{lesson:UR-C129-kalai}



\begin{quote}
Before the new one: say the Urdu for a cheek, then the Urdu for the chin.
\end{quote}

\subsection*{You'll want to know: \ur{کلائی}}
\textbf{\ur{کلائی}} — \emph{kalāī} — \textquotedblleft{}a wrist\textquotedblright{}.

\begin{cousinweb}
\textbf{\ur{کلائی}} (\emph{kalāī}) is a tadbhava, a word inherited from Sanskrit and reshaped by sound change. It comes from \emph{kalācī}, \textquotedblleft{}forearm\textquotedblright{} (CDIAL 2930). The Sanskrit word is known only from the old word-lists.
\end{cousinweb}

\begin{grammarlens}[title={What you've built}]
One more word for this chapter.
\end{grammarlens}

\subsection*{Your turn}
\begin{itemize}
\raggedright
  \item \emph{Say it:} \emph{kalāī}
  \item \emph{Say it:} \emph{kalāī}, once more
  \item \emph{Say it:} say \emph{kalāī}
  \item \emph{Recall:} say the Urdu for a cheek, then the Urdu for the chin, then say \emph{kalāī} again
\end{itemize}

\begin{tcolorbox}[breakable,colback=teal!4,colframe=teal!35!black,title={Before you move on}]
What does \ur{کلائی} mean? (\textquotedblleft{}A wrist\textquotedblright{}.) Say it once more.
\end{tcolorbox}
